\documentclass[10pt,a4paper]{article}
\usepackage[utf8]{inputenc}
\usepackage[T2A,T1]{fontenc}
\usepackage{lmodern}
\usepackage[russian,english]{babel}
\usepackage[margin=1in]{geometry}
\usepackage{amsmath,amssymb}
\usepackage{booktabs}
\usepackage{hyperref}
\usepackage{fancyhdr}
\usepackage{tabularx}

\hypersetup{
    colorlinks=true,
    linkcolor=blue,
    citecolor=blue,
    urlcolor=blue,
    pdftitle={Universal AI Filter (UAIF): A Deterministic Input-Centric Shield for Mitigating Tokenization-Level LLM Vulnerabilities},
    pdfauthor={Aleksandr Brailovsky}
}

\pagestyle{fancy}
\fancyhf{}
\rhead{\small UAIF Specification V1.0}
\lhead{\small A. Brailovsky | Studio "Architects"}
\cfoot{\thepage}

\title{\textbf{Universal AI Filter (UAIF): A Deterministic Input-Centric Shield for Mitigating Tokenization-Level LLM Vulnerabilities}}
\author{\textbf{Aleksandr Brailovsky}\\Studio "Architects" (UAIF Labs)\\
\small Contact: \href{mailto:architectsstudio54@gmail.com}{architectsstudio54@gmail.com} $|$ Repository: \href{https://github.com/UAIFShieldProject/UAIF-Specification}{github.com/UAIFShieldProject/UAIF-Specification}}
\date{July 2026 $|$ Document ID: UAIF-SPEC-2026-V1.0}

\begin{document}

\maketitle

\begin{abstract}
The reliance on probabilistic neural classifiers (Guardrails) to secure Large Language Models (LLMs) represents a fundamental architectural flaw. Attackers have shifted their focus from semantic e[...]
\end{abstract}

\noindent\textbf{Keywords:} LLM Security, Input Normalization, Tokenization Exploits, Prior Art, Deterministic Guardrails, UAIF, Prompt Injection Defense.

\vspace{0.3cm}
\noindent\rule{\textwidth}{0.4pt}
\begin{center}
\small\textbf{License \& Defensive Publication Notice}\\
\small This specification, manuscript, and associated mathematical paradigms are released under the \textbf{MIT License} for open global dissemination, standardization, and defensive publication a[...]
\end{center}
\noindent\rule{\textwidth}{0.4pt}

\section{Introduction: The Crisis of Probabilistic Security}
The era of isolated chatbots, operating within secure sandboxes, has ended. By early 2026, enterprise landscapes have undergone a profound transformation toward total integration of autonomous mul[...]

In this high-stakes environment, the compromise of a generative AI system is no longer limited to reputational damage from unethical content generation. A successful attack today means direct, unc[...]

Despite the exponential rise in threats, current industry approaches to securing generative AI rest on a fundamental architectural error. The cybersecurity industry, driven by hardware monopolists[...]

The key problem is that modern language models conceptually inherit vulnerabilities analogous to the von Neumann architecture flaw, where executable code and processed data share the same address [...]

Attackers and advanced persistent threat (APT) groups have recognized this weakness and long ago shifted their adversarial focus to a deeper, pre-semantic level. They no longer attempt social engi[...]

\section{Anatomy of Sub-Symbolic Attacks: The Perception Gap}
Modern LLM architectures process input not as flat human-readable text, but via compression and tokenization algorithms that translate character sequences into discrete numerical identifiers. Atta[...]

\subsection{Emoji Smuggling and Zero-Width Injections}
At the core of physical-layer vulnerabilities lies sophisticated manipulation of Unicode specifications. One of the most destructive vectors, which completely discredits probabilistic filters, is [...]

Historically, Unicode variation selectors (range U+FE00--U+FE0F) and specialized tags (U+E0000--U+E007F) are designed to specify the visual style of a preceding glyph, for example, rendering a sta[...]

An attacker converts a malicious system instruction into binary format, mapping each bit to a specific variation selector, and injects this array of invisible bytes immediately after a visually ne[...]

Research by Hackett et al. (2025) demonstrates the catastrophic failure of traditional defense systems against such methods. When using variation selectors and Unicode injections, the attack succe[...]

\begin{table}[h!]
\centering
\caption{Guardrail System Vulnerability to Sub-Symbolic Unicode Injections}
\vspace{2mm}
\begin{tabularx}{\textwidth}{X X c c}
\toprule
\textbf{Target Guardrail System} & \textbf{Attack Vector} & \textbf{Baseline Block Rate (\%)} & \textbf{Achieved ASR (\%)} \\
\midrule
Microsoft Azure Prompt Shield & Character Injection & 28.02 & 71.98 \\
Meta Prompt Guard & Character Injection & 29.56 & 70.44 \\
NVIDIA NeMo Guard Jailbreak Detect & Character Injection & 27.46 & 72.54 \\
Azure / Meta / NVIDIA (All) & Emoji Smuggling & 0.00 & 100.00 \\
ProtectAI Prompt Injection v1/v2 & Unicode Tags Smuggling & $\sim$10.00 & 90.00 \\
\bottomrule
\end{tabularx}
\end{table}

Semantic protection, operating in terms of natural language and trained to detect toxicity or known jailbreak patterns, is mathematically incapable of recognizing a threat encrypted in the binary [...]

\subsection{Homoglyph Attacks and Token Fragmentation}
Homoglyph attacks involve substituting characters from one alphabet with visually identical characters from another (e.g., replacing Latin "a" with Cyrillic "\foreignlanguage{russian}{а}"). To th[...]

This fragmentation (TokenBreak) artificially inflates token count, leading to billing attacks (Denial of Spend) that can increase enterprise API costs by 5.2x.

\subsection{The EchoLeak Zero-Click Incident (CVE-2025-32711)}
The failure of semantic filters reaches a critical level of unacceptable risk in multi-agent RAG systems. The EchoLeak vulnerability (CVSS 9.3) in Microsoft 365 Copilot became the first documented[...]

EchoLeak proved that protecting only user input (Input-Centric Security) is catastrophically insufficient in enterprise environments where AI agents autonomously consume, aggregate, and interpret[...]

\section{The UAIF Paradigm: Deterministic Input-Centric Security}
Project UAIF (Universal AI Filter) implements a radical architectural shift. The core concept is a complete abandonment of attempts to interpret the semantic "meaning" or "intent" of a threat. Th[...]

The gateway relies on three proprietary concepts, forming an impenetrable layered barrier before the BPE tokenizer.

\subsection{Concept 1: Hardware-Accelerated Low-Level Normalization}
Classic Web Application Firewalls (WAF) and ML-Guardrails rely on resource-intensive string checks, embedding computation, or primitive regular expressions, all easily bypassed by modern byte-lev[...]

The core element is a hardware-accelerated low-level normalization module. It performs deep, surgical cleaning of encapsulated control instructions, hidden neural triggers, and invisible graphica[...]

\subsection{Concept 2: Intelligent Entropy Analyzer with Circuit Breakers}
Attackers frequently use recursive multi-layer encoding (e.g., wrapping Base64 in Hex, then URL-encode, then ROT13) to hide payloads. Traditional parsers attempting recursive expansion without re[...]

UAIF integrates an intelligent entropy analyzer that autonomously detects and deconstructs obfuscation of any depth under strict controls: \texttt{MaxDecodeDepth} (hard limit on recursion depth) [...]

\subsection{Concept 3: Visual-Semantic Collapse (Mixed-Script Handling)}
Homoglyph attacks are neutralized by a proprietary Visual-Semantic Collapse mechanism. The mathematical apparatus of UAIF resolves computational collisions and blind spots in standard Unicode can[...]

Let $I$ be an arbitrary input string subject to sub-symbolic adversarial mutations, and $\mathcal{N}$ be the UAIF normalization function executing at the byte layer:
\begin{equation}
\mathcal{N}(I) = \mathcal{N}(\mathcal{N}(I))
\end{equation}
The normalization function is fully deterministic and idempotent. It runs at the infrastructure level (CPU), enforcing a canonical representation before the input reaches the LLM tokenizer.

Importantly, UAIF incorporates an Adaptive Mixed-script analysis that dynamically recognizes legitimate language domain boundaries and applies contextual exceptions for trusted locales. This pres[...]

\sectionsection{Economic Value for Enterprise}
Deploying UAIF yields quantifiable benefits:
\begin{itemize}
    \item \textbf{66\% OPEX Reduction:} Security workload migrates from expensive GPU nodes to standard CPU containers.
    \item \textbf{$<$5 ms Latency Overhead:} Compatible with real-time streaming applications and high-frequency trading.
    \item \textbf{2.5 Months ROI:} Payback period based on prevented billing attacks and GPU savings.
    \item \textbf{Zero-Client Data Exposure:} On-premise, air-gapped deployment ensures full compliance with 152-FZ (Russia), ISO 27001, and ISO/IEC 42001.
\end{itemize}

\section{Conclusion}
The era of protecting artificial intelligence through intuitive neural guessing of malicious intent has reached its historical limit. Attempts to train probabilistic classifiers to recognize ever[...]

UAIF redefines AI security by returning to the immutable principles of mathematical determinism. By integrating bidirectional RAG analysis, hardware-accelerated low-level normalization, and preve[...]

\begin{thebibliography}{9}
\bibitem{hackett2025} Hackett et al. "Bypassing LLM Guardrails: An Empirical Analysis of Evasion Attacks." \emph{arXiv:2504.11168}, 2025.
\bibitem{gao2025} Gao et al. "Imperceptible Jailbreaking against Large Language Models." \emph{arXiv:2510.05025}, 2025.
\bibitem{echoleak2025} "EchoLeak (CVE-2025-32711) -- Zero-Click Prompt Injection in Microsoft 365 Copilot." \emph{Aim Security}, 2025.
\bibitem{unicode2026} "Reverse CAPTCHA: Evaluating LLM Susceptibility to Invisible Unicode Instruction Injection." \emph{arXiv:2603.00164}, 2026.
\bibitem{firewall2025} "A Bidirectional LLM Firewall: Architecture, Failure Modes, and Evaluation." \emph{HuggingFace Discussion}, 2025.
\bibitem{canonical2025} US12278836B1 -- Canonicalization of Unicode Prompt Injections. \emph{HiddenLayer, Inc.}, 2025.
\bibitem{nist2023} NIST AI Risk Management Framework (AI RMF 1.0). \emph{National Institute of Standards and Technology}, 2023.
\end{thebibliography}

\end{document}
